\documentclass[12pt]{article}
\usepackage[utf8]{inputenc}
\usepackage[T1]{fontenc}
\usepackage{hyperref}
\title{Semelai Language Resources for K-12 Curriculum Development}
\author{Compiled by Claude Sonnet 5.5 (claude-sonnet-5-5) for Andr\'as Kornai}
\date{}
\begin{document}
\maketitle

\section{Introduction}
Semelai is an Aslian (Austroasiatic, Mon-Khmer) language of the Malay Peninsula, with about 4,100 speakers in Malaysia (2009 figure, per Wikipedia). The best resource found is Nicole Kruspe's reference grammar. No dictionary or school material turned up.

\section{Primary Educational Resources}
No primers, textbooks or school programmes for Semelai were found. The language is not known to be taught in schools.

\section{Linguistic and Reference Materials}
\begin{itemize}
\item Nicole Kruspe, \emph{A Grammar of Semelai}, Cambridge Grammatical Descriptions (2004; paperback 2011). It is the first comprehensive grammar of an Aslian language and covers phonology, morphology, word classes, the verb, pronouns, the noun phrase and syntax. It includes an appendix of texts. \url{https://www.cambridge.org/core/product/0886C124391546CC261AF3BCA4535A5E}
\item Cambridge chapter page for the grammar. \url{https://www.cambridge.org/core/books/grammar-of-semelai/semelai/FEA4C80456E38AD503FA5D5C9B9583AE}
\item LINGUIST List review of the grammar. \url{https://linguistlist.org/issues/18/988/}
\item Wikipedia, Semelai language. \url{https://en.wikipedia.org/wiki/Semelai_language}
\item WALS entry for Semelai. \url{https://wals.info/languoid/lect/wals_code_sml}
\end{itemize}

\section{Religious and Community Resources}
No religious or community-organisation resources were found.

\section{Digital and Audio Resources}
No audio, video or corpus resources were found in the searches.

\section{Curriculum Development Context}
Semelai is used by all in its community, which favours native-language instruction. Published materials are limited to the academic grammar. A curriculum would need new teaching materials and a practical orthography. It could draw on Kruspe's grammar and its texts, with community participation.
\end{document}